problem:  
**Problem (Maximal symmetry dimension for nonflat torsion‑free \(G_2\) metrics with full holonomy).**  
Let \((M^7,g,\varphi)\) be a connected Riemannian manifold with a torsion‑free \(G_2\)-structure, so  
\[
d\varphi = 0, \qquad d{*}\varphi = 0,
\]
and the Levi‑Civita holonomy satisfies \(\mathrm{Hol}(g) = G_2\).  
Let \(\mathfrak{isom}(M,g)\) denote the Lie algebra of local Killing fields of \(g\).

**Open problem.**  
Determine the maximal possible dimension of \(\mathfrak{isom}(M,g)\) among all **nonflat** torsion‑free \(G_2\) manifolds with **full \(G_2\) holonomy**.

In particular:

1. Does there exist a nonflat torsion‑free \(G_2\) metric with \(\dim \mathfrak{isom}(M,g)\ge 3\)?  
   (All known complete examples, such as the Bryant–Salamon metrics, have small isometry algebras of dimension \(2\) or \(3\), but no classifications ensure they are maximal.)

2. More generally, is there a universal upper bound \(D\) such that  
   \[
   \dim \mathfrak{isom}(M,g) \le D
   \]
   for all nonflat, full‑holonomy \(G_2\) manifolds?  
   If so, does equality of this bound characterize a special geometric construction (e.g., cohomogeneity‑one metrics)?

Why is it a "good" problem:  
1. **Precisely circumscribed and avoids settled cases.**  
   Known theorems classify homogeneous Ricci‑flat or locally homogeneous metrics as flat, leaving open only the **partially symmetric** regime between full symmetry and complete rigidity. This problem targets that uncharted margin: exceptional holonomy metrics that admit some isometries but are far from homogeneous.

2. **Distinguishes full holonomy from degenerations.**  
   By restricting to \(\mathrm{Hol}(g)=G_2\), it rules out reducible or product structures (e.g. \(SU(3)\) or \(SU(2)\times SU(2)\) holonomy). Thus it isolates the genuinely exceptional geometry and avoids ambiguity from known lower‑holonomy examples.

3. **Well-defined invariant and quantitative flavor.**  
   The dimension of the Killing algebra is a canonical, algebraic measure of symmetry compatible with local and global analysis, making the question concrete and testable: find or rule out full‑holonomy \(G_2\) metrics with specified symmetry dimension.

4. **Connects to several active research directions.**  
   - The study of cohomogeneity‑one and two \(G_2\) metrics (Bryant–Salamon; Cvetič–Gibbons–Page–Pope) suggests correlations between holonomy reduction and symmetry rank.  
   - The question therefore bridges **representation theory of Lie actions** and **analytic holonomy conditions** from differential geometry, motivating new structure and uniqueness results.

5. **Conceptual simplicity with deep reach.**  
   One sentence captures it—*“What is the maximal symmetry dimension of a nonflat, full holonomy \(G_2\) metric?”*—yet answering it would require synthesizing techniques from isometric group actions, holonomy theory, and analysis of the torsion‑free \(G_2\) PDE. The answer would crystallize how much symmetry exceptional holonomy can bear, a fundamental and currently unresolved problem.